
\documentclass{article}
\usepackage{colortbl}
\usepackage{makecell}
\usepackage{multirow}
\usepackage{supertabular}

\begin{document}

\newcounter{utterance}

\centering \large Interaction Transcript for game `mmlu\_pro', experiment `mmlu\_pro\_philosophy', episode 495 with reflexion{-}llama{-}3.1{-}8b{-}instruct.
\vspace{24pt}

{ \footnotesize  \setcounter{utterance}{1}
\setlength{\tabcolsep}{0pt}
\begin{supertabular}{c@{$\;$}|p{.15\linewidth}@{}p{.15\linewidth}p{.15\linewidth}p{.15\linewidth}p{.15\linewidth}p{.15\linewidth}}
    \# & \multicolumn{2}{c}{Player} && \multicolumn{2}{c}{Game Master} \\
    \hline

    \theutterance \stepcounter{utterance}  
    & & & \multicolumn{4}{p{0.6\linewidth}}{
        \cellcolor[rgb]{0.9,0.9,0.9}{
            \makecell[{{p{\linewidth}}}]{
                \texttt{\tiny{[P1$\langle$GM]}}
                \texttt{You have attempted this same task before and did not succeed. These are the plans you wrote after those attempts.} \\
\texttt{Trial \#1: New plan: } \\
\\ 
\texttt{I failed to achieve the target state of the game because I didn't carefully read the questions and their options. In the first question, I correctly identified the formula, but I didn't notice that the question asked for the best symbolization of the sentence, which implies that there might be multiple correct formulas. In the second question, I chose the correct translation, but I didn't consider the possibility that there might be multiple correct translations. In the third question, I didn't even read the options carefully and chose a random answer. } \\
\\ 
\texttt{Next time, I will carefully read each question and its options, and I will make sure to consider all possible correct answers before making a choice. If I'm still unsure, I will ask for clarification or try to eliminate obviously incorrect options.} \\
\\ 
\\ 
\texttt{The following are multiple choice questions (with answers) about philosophy. Answer only with "The answer is (X)" where X is the correct letter choice.} \\
\\ 
\texttt{Question:} \\
\texttt{Which of the given formulas of PL is the best symbolization of the following sentence?} \\
\texttt{Turtles live long lives and are happy creatures, unless they are injured.} \\
\texttt{Options:} \\
\texttt{A. (L • H) ≡ I} \\
\texttt{B. (L • H) ∨ I} \\
\texttt{C. L • (H ∨ I)} \\
\texttt{D. L • (H ⊃ R).} \\
\texttt{E. N/A} \\
\texttt{F. N/A} \\
\texttt{G. N/A} \\
\texttt{H. N/A} \\
\texttt{I. N/A} \\
\texttt{J. N/A} \\
\texttt{The answer is (B)} \\
\\ 
\texttt{Question:} \\
\texttt{Select the best translation into predicate logic.George borrows Hector's lawnmower. (g: George; h: Hector; l: Hector's lawnmower; Bxyx: x borrows y from z).} \\
\texttt{Options:} \\
\texttt{A. Bhgh} \\
\texttt{B. Bggh} \\
\texttt{C. Bhlh} \\
\texttt{D. Bghl} \\
\texttt{E. Bhlg} \\
\texttt{F. Blhg} \\
\texttt{G. Bllg} \\
\texttt{H. Blgh} \\
\texttt{I. Bhhg} \\
\texttt{J. Bglh} \\
\texttt{The answer is (J)} \\
\\ 
\texttt{Question:} \\
\texttt{The study of reality in the broadest sense, an inquiry into the elemental nature of the universe and the things in it, is known as \_\_\_\_\_.} \\
\texttt{Options:} \\
\texttt{A. quantum physics} \\
\texttt{B. metaphysics} \\
\texttt{C. ontology} \\
\texttt{D. phenomenology} \\
\texttt{E. cosmology} \\
\texttt{F. teleology} \\
\texttt{G. epistemology} \\
\texttt{H. axiology} \\
\texttt{I. N/A} \\
\texttt{J. N/A} \\
\texttt{The answer is (B)} \\
\\ 
\texttt{Question:} \\
\texttt{According to Moore’s “ideal utilitarianism,” the right action is the one that brings about the greatest amount of:} \\
\texttt{Options:} \\
\texttt{A. wealth.} \\
\texttt{B. virtue.} \\
\texttt{C. fairness.} \\
\texttt{D. pleasure.} \\
\texttt{E. peace.} \\
\texttt{F. justice.} \\
\texttt{G. happiness.} \\
\texttt{H. power.} \\
\texttt{I. good.} \\
\texttt{J. knowledge.} \\
\texttt{The answer is (I)} \\
\\ 
\texttt{Question:} \\
\texttt{What is the sign of the covenant for Jewish males?} \\
\texttt{Options:} \\
\texttt{A. Fasting on Yom Kippur} \\
\texttt{B. Lighting Shabbat candles} \\
\texttt{C. The rainbow} \\
\texttt{D. Circumcision} \\
\texttt{E. The Torah} \\
\texttt{F. Bar mitzvah} \\
\texttt{G. Keeping kosher} \\
\texttt{H. Wearing a kippah} \\
\texttt{I. A son} \\
\texttt{J. The Star of David} \\
\texttt{The answer is (D)} \\
\\ 
\texttt{Question:} \\
\texttt{Of the two versions of the principle that Singer considers:} \\
\texttt{Options:} \\
\texttt{A. he considers the stronger version to be correct.} \\
\texttt{B. he considers the weaker version to be correct.} \\
\texttt{C. he proposes a third, completely different principle.} \\
\texttt{D. he argues that the weaker version is stronger.} \\
\texttt{E. he claims that they are equivalent.} \\
\texttt{F. he supports both the stronger and the weaker versions.} \\
\texttt{G. he suggests that the principles are contradictory.} \\
\texttt{H. he does not distinguish between the two principles.} \\
\texttt{I. he states that both principles are incorrect.} \\
\texttt{J. he ends up rejecting them both.} \\
            }
        }
    }
     \\ \\

    \theutterance \stepcounter{utterance}  
    & \multicolumn{4}{p{0.6\linewidth}}{
        \cellcolor[rgb]{0.561,0.737,0.561}{
            \makecell[{{p{\linewidth}}}]{
                \texttt{\tiny{[P1$\rangle$GM]}}
                \texttt{The answer is (B)} \\
\\ 
\texttt{The answer is (J)} \\
\\ 
\texttt{The answer is (B)} \\
\\ 
\texttt{The answer is (I)} \\
\\ 
\texttt{The answer is (D)} \\
\\ 
\texttt{The answer is (B)} \\
            }
        }
    }
    & & \\ \\

    \theutterance \stepcounter{utterance}  
    & & & \multicolumn{2}{p{0.3\linewidth}}{
        \cellcolor[rgb]{0.9,0.9,0.9}{
            \makecell[{{p{\linewidth}}}]{
                \texttt{\tiny{[GM$|$GM]}}
                \texttt{ParseError: Response contains more than a one answer, but the expected format is to only give one.} \\
            }
        }
    }
    & & \\ \\

    \theutterance \stepcounter{utterance}  
    & & & \multicolumn{2}{p{0.3\linewidth}}{
        \cellcolor[rgb]{0.9,0.9,0.9}{
            \makecell[{{p{\linewidth}}}]{
                \texttt{\tiny{[GM$|$GM]}}
                \texttt{game\_result = ABORT} \\
            }
        }
    }
    & & \\ \\

    \theutterance \stepcounter{utterance}  
    & & \multicolumn{4}{p{0.6\linewidth}}{
        \cellcolor[rgb]{0.95,0.95,0.95}{
            \makecell[{{p{\linewidth}}}]{
                \texttt{\tiny{[P1$\leftrightarrow$P]}}
                \texttt{Post{-}episode Reflexion (not sent to GM):} \\
\texttt{New plan: } \\
\\ 
\texttt{I failed to achieve the target state of the game because I didn't carefully read the questions and their options, and I also didn't follow the specified format for answering the questions. In the first question, I correctly identified the formula, but I didn't notice that the question asked for the best symbolization of the sentence, which implies that there might be multiple correct formulas. In the second question, I chose the correct translation, but I didn't consider the possibility that there might be multiple correct translations. In the third question, I didn't even read the options carefully and chose a random answer. } \\
\\ 
\texttt{In the last attempt, I also failed to follow the format for answering the questions, which resulted in the game master rejecting my responses. } \\
\\ 
\texttt{Next time, I will carefully read each question and its options, and I will make sure to consider all possible correct answers before making a choice. If I'm still unsure, I will ask for clarification or try to eliminate obviously incorrect options. I will also make sure to answer the questions in the correct format, using the phrase "The answer is (X)" where X is the correct letter choice.} \\
            }
        }
    }
    & \\ \\

\end{supertabular}
}

\end{document}
